% Generated draft: requires mathematical and source review.
\documentclass[11pt]{article}
\usepackage[margin=25mm]{geometry}
\usepackage{fontspec}
\setmainfont{Latin Modern Roman}
\usepackage{amsmath,amssymb,mathtools,unicode-math}
\setmathfont{Latin Modern Math}
\usepackage{enumitem,needspace,xurl,hyperref}
\setlength{\parindent}{0pt}
\setlength{\parskip}{5pt}
\setlength{\emergencystretch}{3em}
\newcommand{\N}{\mathbb{N}}
\newcommand{\Z}{\mathbb{Z}}
\newcommand{\Q}{\mathbb{Q}}
\newcommand{\R}{\mathbb{R}}
\newcommand{\C}{\mathbb{C}}
\newcommand{\F}{\mathbb{F}}
\newcommand{\A}{\mathbb{A}}
\newcommand{\PP}{\mathbb{P}}
\newcommand{\E}{\mathbb{E}}
\newcommand{\eps}{\varepsilon}
\newcommand{\dd}{\,\mathrm{d}}
\begin{document}
\section*{UnsolvedMath ID: 30001624}
\textbf{Source problem code:} catalogue entry
\subsection{Definitions and mathematical statement}
For a commutative complex \(C^*\)-algebra \(\mathcal L\), let \(X=\operatorname{Spec}(\mathcal L)\) be its character spectrum: the set of nonzero \(*\)-homomorphisms \(\chi:\mathcal L\to\C\), with the topology of pointwise convergence on \(\mathcal L\). Thus \(X\) is locally compact Hausdorff and \(\mathcal L\cong C_0(X)\), where \(C_0(X)\) denotes the continuous complex-valued functions vanishing at infinity.

The supplied record singles out characters belonging to a class denoted by \(\operatorname{Rep}_0(\mathcal L,\C)\). Write
\[
X_0=\{\chi\in X:\chi\in\operatorname{Rep}_0(\mathcal L,\C)\},
\]
and give \(X_0\) the subspace topology inherited from \(X\). The record does not define \(\operatorname{Rep}_0\) or specify the report's construction of \(\mathcal L\); these are unresolved qualifications essential to identifying the intended subset.

\textbf{Local-compactness question.} For the algebra and distinguished character class intended in the cited report, is \(X_0\) locally compact? Explicitly, is it true that
\[
\forall\chi\in X_0\quad\exists U\subseteq X_0:\quad
\chi\in U,\qquad U\text{ is open in }X_0,\qquad
\overline U^{\,X_0}\text{ is compact}?
\]
Here the closure and compactness are taken in \(X_0\). The supplied record does not determine whether the report intends a particular algebra or a quantified class of algebras, so this question cannot be promoted to an assertion about every commutative \(C^*\)-algebra and every distinguished subset of its spectrum.

\textbf{Extension question.} Does the character-spectrum formalism referred to in the report for \(\R^{(\N)}\) admit an analogous construction for separable nuclear spaces? Algebraically,
\[
\R^{(\N)}=\{(x_n)_{n\geq1}:x_n\in\R,\ x_n=0\text{ for all but finitely many }n\}.
\]
The source record does not specify its topology, the scalar field and precise category of the proposed nuclear spaces, the construction to be extended, or the properties that an extension must preserve. Consequently, the extension question is retained as a research target with unresolved specifications; the supplied evidence does not support a fully quantified existence statement.
\subsection{Short English statement}
The first question asks whether the report's distinguished subset of a commutative \(C^*\)-algebra's character spectrum is locally compact in its inherited topology. The second asks whether the associated formalism for \(\R^{(\N)}\) extends to separable nuclear spaces. The distinguished character class and the requirements for the extension are not defined in the supplied record, so both questions retain explicit unresolved qualifications.
\subsection{Sources}
\begin{itemize}
\item Supplied local UnsolvedMath registry record, permanent ID \(30001624\), title \emph{Local Compactness and Nuclear-Space Extensions of Character Spectra}, source problem code \texttt{catalogue entry}, in the Oberwolfach Reports Open Problems collection. No individual proposer is supplied; the proposed year is recorded as 2010.
\item Supplied source citation: \emph{Infinite Dimensional Lie Theory (2011)}, Oberwolfach Reports, pp.\ 2977--3042, DOI \texttt{10.4171/owr/2010/51}. Supplied source URL and literature-source link, labelled ``original OWR report'': \url{https://doi.org/10.4171/owr/2010/51}. The supplied background describes Oberwolfach Reports as documenting mathematical workshops and their open-problem discussions.
\item Literature assessment dated exactly \texttt{2026-08-21}: the supplied curation labels the problem open, describes both questions as research questions in the report, and reports that no general resolution was found. Its listed checked source is the original report linked above. This is supplied dated curation, not an independently verified literature assessment; time-sensitive status and cited claims require independent verification.
\item The registry labels the statement \texttt{corrected\_verified} and supplies a verification note saying that it was checked against the official report and preserves the hypotheses, notation, and open target. These are supplied provenance claims; no independent source consultation or mathematical verification is asserted here.
\end{itemize}
\end{document}
